\chapter{Introduction}
\label{chap-introduction}

\section{Contexte}

Le cours \emph{Initiation aux interfaces graphiques} de Sup Galilée~\cite{toulouse2023}
s'achève par un projet : doter d'une interface graphique, avec la bibliothèque GTK, un jeu
de bataille navale. Le cours ne se réduit pas à l'apprentissage des fenêtres et des boutons.
Il défend une idée d'organisation, la \emph{programmation modulaire}, que ses supports
reprennent d'une séance à l'autre et que l'on peut résumer ainsi :

\begin{itemize}
    \item à chaque concept, un fichier d'en-tête qui en donne la définition et le mode
          d'emploi, et un fichier source qui en réalise les fonctions de manipulation ;
    \item l'utilisation d'un concept est indépendante de sa réalisation ;
    \item un même programme peut recevoir plusieurs interfaces --- console, fenêtre --- et
          c'est alors la seule partie qui doit changer ;
    \item toute fonction écrite est aussitôt testée.
\end{itemize}

De là découle le découpage d'une application en trois bibliothèques : le \emph{modèle}, qui
représente le jeu et ses règles ; la \emph{vue}, qui regroupe les composants graphiques ;
le \emph{contrôleur}, \enquote{chef d'orchestre qui fait le lien entre les deux}.

\section{La première version}

En 2023, nous avons réalisé ce projet en binôme. Plutôt que d'habiller la bibliothèque
fournie, nous avions choisi d'écrire l'intégralité du jeu, pour mieux nous approprier GTK.
Le résultat était un jeu complet et agréable : trois niveaux de difficulté, un placement
aléatoire des flottes, des écrans de règles et de présentation, une identité visuelle
soignée.

Mais ce jeu tenait dans \emph{un seul fichier} de \num{2666} lignes, et son organisation
contredisait ce que le cours enseignait :

\begin{itemize}
    \item l'état de la partie était réparti entre sept variables globales et les composants
          graphiques eux-mêmes : la structure décrivant une case contenait à la fois
          \enquote{y a-t-il un bateau ?} et l'adresse du bouton qui l'affiche ;
    \item les trois niveaux étaient trois copies d'une même fonction de plusieurs centaines de
          lignes ; une correction faite dans l'une manquait dans les autres ;
    \item aucune règle ne pouvait s'exécuter sans ouvrir une fenêtre, donc aucun test
          automatique n'était possible ;
    \item le programme ne fonctionnait que lancé depuis son propre dossier, et sa variante
          sonore que sur la machine où elle avait été écrite.
\end{itemize}

Le programme \emph{marchait}. Il n'était ni vérifiable, ni modifiable sans risque, ni
réutilisable.

\section{Objectifs de la refonte}

La version~2 a un objectif simple à énoncer : \textbf{garder le jeu, refaire tout le reste}.
Le joueur doit retrouver les mêmes règles, les mêmes écrans, les mêmes textes. Sous la
surface, le programme doit devenir ce qu'un logiciel entretenu dans la durée doit être :

\begin{enumerate}
    \item \textbf{modulaire} --- un modèle réellement indépendant de l'interface, au point de
          pouvoir être piloté par une seconde interface ;
    \item \textbf{vérifié} --- des tests automatiques pour les règles, les stratégies de
          l'ordinateur et l'interface elle-même ;
    \item \textbf{reproductible} --- une partie entièrement déterminée par une graine, une
          compilation identique sur un poste de travail et sur un serveur ;
    \item \textbf{livrable} --- un exécutable unique, des paquets, une intégration et une
          livraison continues.
\end{enumerate}

\section{Organisation du rapport}

Le \cref{chap-exigences} fixe les règles du jeu et les exigences. Le
\cref{chap-architecture} présente l'architecture et la justifie. Les
\cref{chap-modele,chap-adversaire} détaillent le moteur de jeu puis l'adversaire
artificiel, avec les mesures de sa force. Le \cref{chap-interface} décrit l'interface
graphique. Le \cref{chap-qualite} expose la démarche de vérification et la chaîne
d'intégration continue. Le \cref{chap-bilan} dresse le bilan chiffré de la refonte, ses
limites et ses prolongements possibles.
